% Complete source-bound development results; lower MSE is better.
\begin{table}[!htbp]\centering\small
\setlength{\tabcolsep}{4pt}
\begin{tabular}{@{}lrrrr@{}}\toprule
& \multicolumn{2}{c}{Native $4\!\times\!4$ MSE $\downarrow$}
& \multicolumn{2}{c}{Original $2\!\times\!2$ MSE $\downarrow$} \\
Method & $h=5$ & $h=10$ & $h=5$ & $h=10$ \\ \midrule
Autoregressive & 0.149573 & 0.203950 & 0.148577 & 0.202916 \\
Additive anchor & 0.149606 & 0.200244 & 0.148705 & 0.198704 \\
ShiftWM (ours) & \textbf{0.145074} & \textbf{0.193141} & 0.147921 & 0.197766 \\
No context (ours, ablation) & 0.145143 & 0.193330 & \textbf{0.147635} & \textbf{0.197553} \\
No actions (ours, ablation) & 0.149480 & 0.202249 & 0.152414 & 0.206188 \\
DINO-WM: shifted one-step & 0.219267 & 0.355590 & 0.224146 & 0.369692 \\
DINO-WM: recursive H10 & 0.233920 & 0.331548 & 0.238623 & 0.339783 \\
Persistence & 0.164028 & 0.226129 & 0.163156 & 0.225622 \\
\bottomrule\end{tabular}
\caption{Posthoc DROID development comparison with adapted DINO-WM.
Both external training objectives retain all three seeds and 30 epochs:
the official shifted one-step objective and matched recursive ten-step training.
All learned models use the same training/development split; means weight windows
within episode, 141 episodes, and three seeds equally (1,631 windows, 59 sessions).
Black bold marks the lowest point mean, not significance. Native and original
pooled coordinates have separate training scales; their errors are not interchangeable.
$h$ counts action blocks. This is not an official benchmark reproduction or a
control-success comparison.}
\label{tab:external-dinowm-scores}\end{table}
